\section{Translate sixteen sextets with comparison masks}
The vector encoder can apply the alphabet's piecewise offsets without sixteen table
lookups. Start from the sextets and accumulate corrections: add six above 25, subtract
75 above 51, subtract fifteen above 61, add three when equal to 63, then add 65 to every
lane. This gives the five alphabet cases in the earlier table.

A true vector comparison produces an all-one byte, hexadecimal \texttt{ff}; false
produces zero. AND with six changes a true mask into the correction six. Adding the
raw comparison mask would add 255 modulo 256, which is not the required correction.

\Needspace{9\baselineskip}
\MARSCodeLabel{Convert a comparison mask into an arithmetic correction.}
\begin{minted}[samepage=true]{mars-idea}
    ldi.b $w2, 25
    clt_u.b $w3, $w2, $w1
    andi.b $w3, $w3, 6
    addv.b $w4, $w1, $w3
\end{minted}

Retain the original sextets in \texttt{w1} for every threshold test while accumulating
the corrected values in \texttt{w4}. Comparing a partly translated character with a
sextet threshold would mix the two representations.

The byte additions and subtractions wrap modulo 256. Intermediate \texttt{w4} lanes
need not yet be ASCII. For sextet 52, the sequence before the final addition is
\(52+6=58\), then \(58-75\equiv239\pmod{256}\); adding 65 finally produces 48,
ASCII \texttt{0}. For sextet 63, the successive byte values are 69, 250, 235, 238 and
47, ending at \texttt{/}. Inspect the result after all corrections before treating it
as text.

The final \texttt{st.b} writes exactly sixteen bytes to the kernel's destination.
The input contract supplies twelve bytes and the output contract supplies sixteen;
those footprints are separate from the register's fixed width.
